% GENERATED FILE. Edit canonical lessons, then run npm run generate:books.
% canonical-source-hash: fnv1a64:a55804fab19fe28a
% canonical-lessons: ML-C274-niru, ML-C274-ulukku, ML-C274-otivu, ML-C274-pollal, ML-C274-prameham

\chapter{Swelling, Sprain, Fracture, Burn, Diabetes}
\label{ch:ml-swelling-sprain-fracture-burn-diabetes}
\hlchaptermodality{274}

\begin{chapteropening}
\textbf{By the end of this chapter:} \emph{I can say a swelling, a sprain, a fracture, a burn and diabetes.}

Complete the last lesson of chapter 274: I can say a swelling, a sprain, a fracture, a burn and diabetes.
\end{chapteropening}

\section[\texorpdfstring{\ml{നീര്} (nīrŭ)}{നീര് (nīrŭ)}]{\ml{നീര്} (nīrŭ) — a swelling}

\label{lesson:ML-C274-niru}



\begin{quote}
Before the new one: say the Malayalam for to spin, then the Malayalam for to pinch.
\end{quote}

\subsection*{You'll want to know: \ml{നീര്}}
\textbf{\ml{നീര്}} — \emph{nīrŭ} — \textquotedblleft{}a swelling\textquotedblright{}.

\begin{grammarlens}[title={What you've built}]
One more word for this chapter.
\end{grammarlens}

\subsection*{Your turn}
\begin{itemize}
\raggedright
  \item \emph{Say it:} \emph{nīrŭ}
  \item \emph{Say it:} \emph{nīrŭ}, once more
  \item \emph{Say it:} say \emph{nīrŭ}
  \item \emph{Recall:} say the Malayalam for to spin, then the Malayalam for to pinch, then say \emph{nīrŭ} again
\end{itemize}

\begin{tcolorbox}[breakable,colback=teal!4,colframe=teal!35!black,title={Before you move on}]
What does \ml{നീര്} mean? (\textquotedblleft{}A swelling\textquotedblright{}.) Say it once more.
\end{tcolorbox}
\section[\texorpdfstring{\ml{ഉളുക്ക്} (uḷukkŭ)}{ഉളുക്ക് (uḷukkŭ)}]{\ml{ഉളുക്ക്} (uḷukkŭ) — a sprain}

\label{lesson:ML-C274-ulukku}



\begin{quote}
Before the new one: say the Malayalam for to pinch, then the Malayalam for a swelling.
\end{quote}

\subsection*{You'll want to know: \ml{ഉളുക്ക്}}
\textbf{\ml{ഉളുക്ക്}} — \emph{uḷukkŭ} — \textquotedblleft{}a sprain\textquotedblright{}.

\begin{grammarlens}[title={What you've built}]
One more word for this chapter.
\end{grammarlens}

\subsection*{Your turn}
\begin{itemize}
\raggedright
  \item \emph{Say it:} \emph{uḷukkŭ}
  \item \emph{Say it:} \emph{uḷukkŭ}, once more
  \item \emph{Say it:} say \emph{uḷukkŭ}
  \item \emph{Recall:} say the Malayalam for to pinch, then the Malayalam for a swelling, then say \emph{uḷukkŭ} again
\end{itemize}

\begin{tcolorbox}[breakable,colback=teal!4,colframe=teal!35!black,title={Before you move on}]
What does \ml{ഉളുക്ക്} mean? (\textquotedblleft{}A sprain\textquotedblright{}.) Say it once more.
\end{tcolorbox}
\section[\texorpdfstring{\ml{ഒടിവ്} (oṭivŭ)}{ഒടിവ് (oṭivŭ)}]{\ml{ഒടിവ്} (oṭivŭ) — a fracture}

\label{lesson:ML-C274-otivu}



\begin{quote}
Before the new one: say the Malayalam for a swelling, then the Malayalam for a sprain.
\end{quote}

\subsection*{You'll want to know: \ml{ഒടിവ്}}
\textbf{\ml{ഒടിവ്}} — \emph{oṭivŭ} — \textquotedblleft{}a fracture\textquotedblright{}.

\begin{grammarlens}[title={What you've built}]
One more word for this chapter.
\end{grammarlens}

\subsection*{Your turn}
\begin{itemize}
\raggedright
  \item \emph{Say it:} \emph{oṭivŭ}
  \item \emph{Say it:} \emph{oṭivŭ}, once more
  \item \emph{Say it:} say \emph{oṭivŭ}
  \item \emph{Recall:} say the Malayalam for a swelling, then the Malayalam for a sprain, then say \emph{oṭivŭ} again
\end{itemize}

\begin{tcolorbox}[breakable,colback=teal!4,colframe=teal!35!black,title={Before you move on}]
What does \ml{ഒടിവ്} mean? (\textquotedblleft{}A fracture\textquotedblright{}.) Say it once more.
\end{tcolorbox}
\section[\texorpdfstring{\ml{പൊള്ളൽ} (poḷḷal)}{പൊള്ളൽ (poḷḷal)}]{\ml{പൊള്ളൽ} (poḷḷal) — a burn}

\label{lesson:ML-C274-pollal}



\begin{quote}
Before the new one: say the Malayalam for a sprain, then the Malayalam for a fracture.
\end{quote}

\subsection*{You'll want to know: \ml{പൊള്ളൽ}}
\textbf{\ml{പൊള്ളൽ}} — \emph{poḷḷal} — \textquotedblleft{}a burn\textquotedblright{}.

\begin{grammarlens}[title={What you've built}]
One more word for this chapter.
\end{grammarlens}

\subsection*{Your turn}
\begin{itemize}
\raggedright
  \item \emph{Say it:} \emph{poḷḷal}
  \item \emph{Say it:} \emph{poḷḷal}, once more
  \item \emph{Say it:} say \emph{poḷḷal}
  \item \emph{Recall:} say the Malayalam for a sprain, then the Malayalam for a fracture, then say \emph{poḷḷal} again
\end{itemize}

\begin{tcolorbox}[breakable,colback=teal!4,colframe=teal!35!black,title={Before you move on}]
What does \ml{പൊള്ളൽ} mean? (\textquotedblleft{}A burn\textquotedblright{}.) Say it once more.
\end{tcolorbox}
\section[\texorpdfstring{\ml{പ്രമേഹം} (pramēhaṁ)}{പ്രമേഹം (pramēhaṁ)}]{\ml{പ്രമേഹം} (pramēhaṁ) — diabetes}

\label{lesson:ML-C274-prameham}



\begin{quote}
Before the new one: say the Malayalam for a fracture, then the Malayalam for a burn.
\end{quote}

\subsection*{You'll want to know: \ml{പ്രമേഹം}}
\textbf{\ml{പ്രമേഹം}} — \emph{pramēhaṁ} — \textquotedblleft{}diabetes\textquotedblright{}.

\begin{grammarlens}[title={What you've built}]
One more word for this chapter.
\end{grammarlens}

\subsection*{Your turn}
\begin{itemize}
\raggedright
  \item \emph{Say it:} \emph{pramēhaṁ}
  \item \emph{Say it:} \emph{pramēhaṁ}, once more
  \item \emph{Say it:} say \emph{pramēhaṁ}
  \item \emph{Recall:} say the Malayalam for a fracture, then the Malayalam for a burn, then say \emph{pramēhaṁ} again
\end{itemize}

\begin{tcolorbox}[breakable,colback=teal!4,colframe=teal!35!black,title={Before you move on}]
What does \ml{പ്രമേഹം} mean? (\textquotedblleft{}Diabetes\textquotedblright{}.) Say it once more.
\end{tcolorbox}
